% 版式参数。必须在 \documentclass 之前读入。
%
% 四种版式由构建脚本通过 \def\booklayout{...} 指定：
%   screen   A4，屏幕阅读（默认）
%   print    A4，印刷装订，内侧留装订线
%   ereader  6 寸电子书阅读器（91mm × 122mm），灰度，窄边距
%   phone    手机竖屏（88mm × 176mm），彩色，窄边距

\providecommand{\booklayout}{screen}

\newif\ifbknarrow
\newif\ifbkgrayscale
\newif\ifbkprint
\bknarrowfalse
\bkgrayscalefalse
\bkprintfalse

\def\bkScreen{screen}
\def\bkPrint{print}
\def\bkEreader{ereader}
\def\bkPhone{phone}

\ifx\booklayout\bkPrint
  \bkprinttrue
\fi
% 灰度只用于电子墨水屏；手机屏是彩色的，保留彩色版式。
\ifx\booklayout\bkEreader
  \bknarrowtrue
  \bkgrayscaletrue
\fi
\ifx\booklayout\bkPhone
  \bknarrowtrue
\fi

% 字号：标准文档类只接受 10/11/12pt，窄版式先用 10pt 建类，
% 再由 preamble 里的 \AtBeginDocument 把正文缩到 9pt
% （6 寸屏上 10pt 相当于 A4 的 14pt，偏大）。
\ifbknarrow
  \def\bkfontsize{10pt}
\else
  \def\bkfontsize{11pt}
\fi

% 单双面：印刷版双面装订，章与部分从右页起排；其余版式单面，
% 避免 \mainmatter 与 \part 为凑奇数页插入空白页。
\ifbkprint
  \def\bkside{twoside}
\else
  \def\bkside{oneside}
\fi
